\documentclass[12pt]{article}
\usepackage[utf8]{inputenc}
\usepackage[T1]{fontenc}
\usepackage{hyperref}
\title{Susu Language Resources for K-12 Curriculum Development}
\author{Compiled by Claude Sonnet 5.5 (claude-sonnet-5-5) for Andr\'as Kornai}
\date{}
\begin{document}
\maketitle

\section{Introduction}
Susu (Soso, Sosoxui; ISO 639-3 \texttt{sus}) is a Mande language spoken mainly in Guinea, with a community in Sierra Leone. Its closest relative is Yalunka, and it is one of Guinea's national languages. Search results were thin on modern educational material. The earliest printed works are missionary-era.

\section{Primary Educational Resources}
No modern primer or textbook was found in search results. Early printed material, per Wikipedia, includes a spelling book and Anglican catechism (1801) and Henry Brunton's \emph{A grammar and vocabulary of the Susoo language} (1802). A Susu--French lexicon of about 88 pages from Guinea's literacy service appears in a bibliography, but no link was found.

\section{Linguistic and Reference Materials}
\begin{itemize}
\item Wikipedia, ``Susu language'': \url{https://en.wikipedia.org/wiki/Susu_language}
\item Omniglot, Susu (Sosoxui): \url{https://www.omniglot.com/writing/susu.htm}
\item CoCoon (Huma-Num) archive lookup for \texttt{sus}: \url{https://cocoon.huma-num.fr/exist/crdo/lexvo.xql?code=sus}
\item Sierra Leone Heritage glossary entry on Soso: \url{https://sierraleoneheritage.org/glossary/soso}
\end{itemize}

\section{Religious and Community Resources}
\begin{itemize}
\item Scripture Earth listing for \texttt{sus} (Scripture texts and Gospel films): \url{https://scriptureearth.org/00eng.php?iso=sus}
\item YouVersion Susu text (Matthew, 1816 edition): \url{https://www.bible.com/yo/versions/2066}
\end{itemize}

\section{Digital and Audio Resources}
\begin{itemize}
\item Global Recordings Network, Susu (Sousou) audio: \url{https://globalrecordings.net/en/language/sus}
\end{itemize}

\section{Curriculum Development Context}
Susu is widely spoken and recognised as a national language in Guinea, but the reference description says it is not known to be taught in schools, and no school materials were found. Further searching for a current Susu--French or Susu--English dictionary and for Guinean literacy-service publications would be the next step.

\end{document}